\chapter{Análisis de Riesgos}
\label{ch:riesgos}

Este capítulo identifica y analiza los principales riesgos técnicos,
operacionales y de alcance asociados al \textbf{CMOS Microscopy Payload (CMOSM)}.
El objetivo es reconocer de manera explícita las limitaciones del sistema en su
fase actual de desarrollo y documentar las acciones de mitigación consideradas,
siguiendo buenas prácticas de ingeniería de sistemas (ECSS/NASA).

El análisis de riesgos se enfoca en la operación del payload como demostrador
tecnológico en entorno de laboratorio, sin considerar aún condiciones de vuelo.

% -------------------------------------------------------------------
\section{Metodología de análisis de riesgos}

Los riesgos se clasifican de acuerdo con:
\begin{itemize}
    \item \textbf{Tipo}: técnico, operacional o de alcance.
    \item \textbf{Probabilidad}: baja (B), media (M) o alta (A).
    \item \textbf{Impacto}: bajo (B), medio (M) o alto (A).
\end{itemize}

Para cada riesgo identificado se define una estrategia de mitigación o
contención, acorde con el estado actual del proyecto.

% -------------------------------------------------------------------
\section{Riesgos técnicos}

\subsection{Acoplamiento óptico muestra--sensor}

\textbf{Descripción:}  
La calidad de imagen en microscopía de contacto depende críticamente de la
distancia y uniformidad del contacto entre la muestra y el sensor CMOS. Variaciones
mínimas pueden degradar el contraste y la definición de bordes.

\textbf{Probabilidad:} Alta  
\textbf{Impacto:} Alto  

\textbf{Mitigación:}
\begin{itemize}
    \item Uso de procedimientos experimentales controlados para preparación de
    muestras.
    \item Evaluación sistemática de distintas interfaces muestra--sensor.
    \item Registro de condiciones experimentales en los metadatos.
\end{itemize}

\subsection{Estabilidad de iluminación}

\textbf{Descripción:}  
Derivas temporales en la intensidad del OLED pueden afectar la reproducibilidad
entre capturas consecutivas.

\textbf{Probabilidad:} Media  
\textbf{Impacto:} Medio  

\textbf{Mitigación:}
\begin{itemize}
    \item Calentamiento previo del sistema antes de captura.
    \item Monitoreo de métricas de histograma y saturación.
    \item Uso de patrones uniformes en operación nominal.
\end{itemize}

\subsection{Ruido del sensor CMOS}

\textbf{Descripción:}  
El ruido térmico y el ruido de patrón fijo del sensor pueden limitar la calidad de
imagen en exposiciones largas.

\textbf{Probabilidad:} Media  
\textbf{Impacto:} Medio  

\textbf{Mitigación:}
\begin{itemize}
    \item Aplicación de calibración mediante dark frame y flat field.
    \item Operación a temperaturas controladas de laboratorio.
\end{itemize}

\subsection{Limitaciones de reconstrucción FPM}

\textbf{Descripción:}  
Las técnicas de reconstrucción FPM pueden no converger adecuadamente debido a
limitaciones de estabilidad mecánica, alineamiento o acoplamiento óptico.

\textbf{Probabilidad:} Alta  
\textbf{Impacto:} Bajo  

\textbf{Mitigación:}
\begin{itemize}
    \item Clasificación de FPM como modo experimental.
    \item No dependencia de FPM para operación nominal del payload.
\end{itemize}

% -------------------------------------------------------------------
\section{Riesgos operacionales}

\subsection{Preparación y manipulación de muestras}

\textbf{Descripción:}  
La manipulación manual de muestras biológicas puede introducir variabilidad entre
ensayos y riesgo de contaminación del sensor.

\textbf{Probabilidad:} Media  
\textbf{Impacto:} Medio  

\textbf{Mitigación:}
\begin{itemize}
    \item Procedimientos estandarizados de preparación.
    \item Limpieza controlada del sensor.
    \item Uso de etiquetas experimentales en metadatos.
\end{itemize}

\subsection{Errores de configuración de captura}

\textbf{Descripción:}  
Configuraciones incorrectas de exposición o ganancia pueden provocar saturación o
subexposición.

\textbf{Probabilidad:} Baja  
\textbf{Impacto:} Medio  

\textbf{Mitigación:}
\begin{itemize}
    \item Vista previa en tiempo real con histograma.
    \item Bloqueo de parámetros automáticos del sensor.
\end{itemize}

% -------------------------------------------------------------------
\section{Riesgos de alcance y proyecto}

\subsection{Extensión del alcance del procesamiento computacional}

\textbf{Descripción:}  
Existe el riesgo de sobreextender el alcance del proyecto hacia técnicas
computacionales avanzadas que no aporten mejoras significativas a la imagen base.

\textbf{Probabilidad:} Media  
\textbf{Impacto:} Bajo  

\textbf{Mitigación:}
\begin{itemize}
    \item Definición clara del modo nominal de operación.
    \item Evaluación objetiva del aporte incremental de cada técnica.
\end{itemize}

\subsection{Transición a entorno espacial}

\textbf{Descripción:}  
El sistema no ha sido aún verificado para condiciones de vibración, vacío o
radiación propias del entorno espacial.

\textbf{Probabilidad:} Alta  
\textbf{Impacto:} Alto  

\textbf{Mitigación:}
\begin{itemize}
    \item Delimitación explícita del alcance del presente CDR a entorno de
    laboratorio.
    \item Identificación de ensayos ambientales como trabajo futuro.
\end{itemize}

% -------------------------------------------------------------------
\section{Resumen de riesgos}

\begin{table}[H]
\centering
\begin{tabular}{p{4cm} p{3cm} p{3cm} p{3cm}}
\toprule
\textbf{Riesgo} & \textbf{Tipo} & \textbf{Probabilidad} & \textbf{Impacto} \\
\midrule
Acoplamiento muestra--sensor & Técnico & Alta & Alto \\
Estabilidad de iluminación & Técnico & Media & Medio \\
Ruido del sensor & Técnico & Media & Medio \\
Reconstrucción FPM & Técnico & Alta & Bajo \\
Preparación de muestras & Operacional & Media & Medio \\
Errores de configuración & Operacional & Baja & Medio \\
Transición a vuelo & Alcance & Alta & Alto \\
\bottomrule
\end{tabular}
\caption{Resumen de riesgos identificados para el CMOSM.}
\end{table}

% -------------------------------------------------------------------
\section{Conclusión del análisis de riesgos}

El análisis de riesgos realizado permite identificar claramente los factores
críticos que afectan el desempeño del CMOSM en su fase actual. Las mitigaciones
propuestas son coherentes con el carácter experimental del payload y establecen
una base sólida para futuras iteraciones del proyecto INTISAT.
